\documentclass[11pt, twocolumn, landscape, a4paper]{article}

\usepackage[margin=1cm]{geometry}
\usepackage[utf8]{inputenc}
\usepackage[spanish]{babel}
\usepackage{amssymb, amsmath, amsbsy}
\usepackage{mathpazo}
\usepackage{tasks}
\usepackage{verbatim}
\usepackage{graphicx} 
\usepackage[pdftex]{hyperref}
\hypersetup{pdftitle={Ejercicios para Funciones en álgebra},pdfauthor={Luis Carrillo Gutiérrez}}

\renewcommand{\familydefault}{\sfdefault}

\begin{document}
\subsection*{Funciones}

\begin{itemize}
\item{En los siguientes conjuntos. Identificar ¿cuáles cumplen como \emph{funciones}? :
\begin{itemize}
\item{$F = \Big\{ (4;\,5),\;(2;\,3),\;(7;\,9)\Big\}$}
\item{$G = \Big\{ (3;\,8),\;(3;\,2),\;(7;\,4)\Big\}$}
\item{$H = \Big\{ (2;\,9),\;(3;\,8),\;(2;\,9),\;(7;\,6)\Big\}$}
\end{itemize}

\begin{tasks}(5)
\task{$F$ y $H$}\task{$F$ y $G$}\task{$G$}\task{$G$ y $H$}\task{$H$}
\end{tasks}

\begin{comment}
Nos piden determinar cuáles de los siguientes conjuntos son funciones: F = {(4;5), (2;3), (7;9)} Las primeras componentes son diferentes, entonces f es una _función_. G = {(3:6),(3;2),(7;4)}. Se repiten las primeras componentes, entonces g NO es una _función_. H = {(2;9),(3;6),(2;9),(7;6)} 
Pares ordenados iguales, se asumen como uno solo. Luego en H = {(2;9),(3;8),(7;6)}. Las primeras componentes son diferentes, entonces h es función. Por lo tanto, los conjuntos *F* y *H* son funciones.
\end{comment}
}
\item{Si \scalebox{1.2}{$f$} es una función (expresada como conjunto de pares ordenados) : \\ $f = \Big\{(2;\,a + 1),\;(3;\,8),\;(2;\,11),\;(3;\,3b - 4),\;(7;\,11) \Big\}$ \\ 
Hallar el valor de :\quad\scalebox{1.25}{$a + b$}
\begin{tasks}(5)
\task{10}\task{12}\task{13}\task{14}\task{21}
\end{tasks}
}
\item{Hallar/Determinar el {\tt dominio} de la función : $f(x) = \sqrt{\,3x-9\,}$
\begin{tasks}(3) 
\task{$\langle\,-\infty\,;\,-3\,\rangle$}
\task{$\langle\,-\infty\,;\,3\,\rangle$}
\task{$\langle\,3\,;\,+\infty\,\rangle$}
\task{$\langle\,-\infty\,;\,9\,\rangle$}
\task{$[\,9\,;\,\infty\,\rangle$}
\end{tasks}
}
\item{Hallar/Determinar el {\tt dominio} de la función : $f(x) = \sqrt{\,x + 2\;} + \sqrt{\,8 - x\;}$
\begin{tasks}(5) 
\task{$\big[\,2\,;\,8\,\big]$}
\end{tasks}
}
\item{Hallar/Determinar el {\tt dominio} de la función : $f(x) = \dfrac{x - 4}{\;x - 2\;}$
\begin{tasks}(5)
\task{$\langle\,2\,;\,4\,\rangle$}
\task{$[\,2\,;\,4\,]$} 
\task{$\langle\,2\,;\,4\,]$} 
\task{$\langle\,2\,;\,+\infty\,\rangle$} 
\task{$\mathbb{R} - \{2\}$}
\end{tasks}
}
\item{Hallar/Determinar el {\tt dominio} de la función : $f(x) = \dfrac{2}{\;x^{2} - 4\;}$
\begin{tasks}(3) 
\task{$\mathbb{R}$}\task{$\{2\}$}\task{$\{-2\}$}
\task{$\mathbb{R} - \{2\}$}\task{$\mathbb{R} - \{2\,;\,-2\}$}
\end{tasks}
}
\item{Calcular el {\tt rango} de : $f(x) = 3x + 2$;\quad si $x \in \big\langle\,2\,;\,5\,\big]$
\begin{tasks}(5)
\task{$\langle\,0\,;\,5\,\rangle$} 
\task{$\langle\,3\,;\,10\,\rangle$} 
\task{$\langle\,8\,;\,17\,\rangle$}
\task{$\langle\,18\,;\,29\,\rangle$}
\task{$\langle\,2\,;\,3\,\rangle$}
\end{tasks}
}
\item{Si $x \in \langle-2,\,-1\rangle$.\quad Calcular el {\tt rango} de la función \scalebox{1.25}{$f$}, cuya regla de correspondencia es : $f(x) = \dfrac{\;1 - x\;}{2}$
\begin{tasks}(3) 
\task{$\Big\langle\,0\,;\,\dfrac{3}{2}\,\Big\rangle$} 
\task{$\Big\langle\,-1\,;\,\dfrac{3}{2}\,\Big\rangle$} 
\task{$\Big\langle\,-1\,;\,\dfrac{1}{2}\,\Big\rangle$} 
\task{$\Big\langle\,1\,;\,\dfrac{3}{2}\,\Big\rangle$}
\task{$\Big\langle\,0\,;\,\dfrac{1}{2}\,\Big\rangle$}
\end{tasks}
}
\item{Calcular el {\tt rango} de : $f(x) = x^{2} + 2x + 3$, donde $x \in \big[\,5\,;\,10\big\rangle$
\begin{tasks}(3) 
\task{$\langle\,10\,;\,51\,\rangle$}
\task{$\langle\,38\,;\,123\,\rangle$}
\task{$\langle\,43\,;\,123\,\rangle$}
\end{tasks}
}
\item{Calcular el {\tt rango} de la función \scalebox{1.25}{$f$} tiene como regla de correspondencia a :\qquad $f(x) = \dfrac{\;x + 2\;}{x + 1}$\quad y como {\tt dominio} al intervalo $\langle\,2\,;\,3\,\rangle$
\begin{tasks}(3) 
\task{$\Big\langle\,0\,;\,\dfrac{4}{3}\,\Big\rangle$}
\task{$\Big\langle\,\dfrac{3}{4}\,;\,\dfrac{4}{3}\,\Big\rangle$}
\task{$\Big\langle\,\dfrac{4}{3}\,;\,\dfrac{5}{4}\,\Big\rangle$}
\task{$\Big\langle\,\dfrac{1}{4}\,;\,\dfrac{1}{2}\,\Big\rangle$}
\task{$\Big\langle\,\dfrac{5}{4}\,;\,\dfrac{4}{3}\,\Big\rangle$}
\end{tasks}
}
\end{itemize}
\end{document}